% GENERATED FILE. Edit canonical lessons, then run npm run generate:books.
% canonical-source-hash: fnv1a64:b61491f9868d048d
% canonical-lessons: BN-C340-boidho, BN-C340-nishiddho, BN-C340-dhara, BN-C340-palon-kora, BN-C340-proyog-kora

\chapter{What the Law Allows}
\label{ch:bn-what-the-law-allows}
\hlchaptermodality{340}

\begin{chapteropening}
\textbf{By the end of this chapter:} \emph{I can say what the law allows, forbids and requires.}

Complete the last lesson of chapter 340: I can say what the law allows, forbids and requires.
\end{chapteropening}

\section[\texorpdfstring{\bn{বৈধ} (boidhô)}{বৈধ (boidhô)}]{\bn{বৈধ} (boidhô) — lawful; valid}

\label{lesson:BN-C340-boidho}



\begin{quote}
Before the new one: say the Bengali for to defend, then the Bengali for to enjoy a right.
\end{quote}

\subsection*{You'll want to know: \bn{বৈধ}}
\textbf{\bn{বৈধ}} — \emph{boidhô} — \textquotedblleft{}lawful; valid\textquotedblright{}.

\begin{quote}
\bn{তোমার} \bn{ভিসা} \bn{কি} \bn{বৈধ}?
\end{quote}

\begin{quote}
\emph{tomār bhisā ki boidhô?}
\end{quote}

\begin{quote}
Is your visa valid?
\end{quote}

\begin{cousinweb}
\textbf{\bn{বৈধ}} (\emph{boidhô}) is Sanskrit \emph{vaidha}, 'according to rule', legal (Monier-Williams), from \emph{vidhi} 'a rule'. A paper or an act can be \bn{বৈধ}.
\end{cousinweb}

\begin{grammarlens}[title={What you've built}]
One more word for giving reasons and opinions.
\end{grammarlens}

\subsection*{Your turn}
\begin{itemize}
\raggedright
  \item \emph{Say it:} \emph{boidhô}
  \item \emph{Say it:} \emph{boidhô}, once more
  \item \emph{Say it:} \emph{tomār bhisā ki boidhô?} — Is your visa valid?
  \item \emph{Recall:} say the Bengali for to defend, then the Bengali for to enjoy a right, then say \emph{boidhô} again
\end{itemize}

\begin{tcolorbox}[breakable,colback=teal!4,colframe=teal!35!black,title={Before you move on}]
What does \bn{বৈধ} mean? (\textquotedblleft{}Lawful; valid\textquotedblright{}.) Say it once more.
\end{tcolorbox}
\section[\texorpdfstring{\bn{নিষিদ্ধ} (nishiddhô)}{নিষিদ্ধ (nishiddhô)}]{\bn{নিষিদ্ধ} (nishiddhô) — forbidden; banned}

\label{lesson:BN-C340-nishiddho}



\begin{quote}
Before the new one: say the Bengali for to enjoy a right, then the Bengali for lawful.
\end{quote}

\subsection*{You'll want to know: \bn{নিষিদ্ধ}}
\textbf{\bn{নিষিদ্ধ}} — \emph{nishiddhô} — \textquotedblleft{}forbidden; banned\textquotedblright{}.

\begin{quote}
\bn{এখানে} \bn{পার্কিং} \bn{নিষিদ্ধ।}
\end{quote}

\begin{quote}
\emph{ekhāne pārking nishiddhô.}
\end{quote}

\begin{quote}
Parking is forbidden here.
\end{quote}

\begin{cousinweb}
\textbf{\bn{নিষিদ্ধ}} (\emph{nishiddhô}) is Sanskrit \emph{niṣiddha}, kept back or forbidden (Monier-Williams). It is the word on notices: \bn{ধূমপান} \bn{নিষিদ্ধ}, no smoking.
\end{cousinweb}

\begin{grammarlens}[title={What you've built}]
One more word for giving reasons and opinions.
\end{grammarlens}

\subsection*{Your turn}
\begin{itemize}
\raggedright
  \item \emph{Say it:} \emph{nishiddhô}
  \item \emph{Say it:} \emph{nishiddhô}, once more
  \item \emph{Say it:} \emph{ekhāne pārking nishiddhô.} — Parking is forbidden here.
  \item \emph{Recall:} say the Bengali for to enjoy a right, then the Bengali for lawful, then say \emph{nishiddhô} again
\end{itemize}

\begin{tcolorbox}[breakable,colback=teal!4,colframe=teal!35!black,title={Before you move on}]
What does \bn{নিষিদ্ধ} mean? (\textquotedblleft{}Forbidden; banned\textquotedblright{}.) Say it once more.
\end{tcolorbox}
\section[\texorpdfstring{\bn{ধারা} (dhārā)}{ধারা (dhārā)}]{\bn{ধারা} (dhārā) — a section (of a law); a clause}

\label{lesson:BN-C340-dhara}



\begin{quote}
Before the new one: say the Bengali for lawful, then the Bengali for banned.
\end{quote}

\subsection*{You'll want to know: \bn{ধারা}}
\textbf{\bn{ধারা}} — \emph{dhārā} — \textquotedblleft{}a section (of a law); a clause\textquotedblright{}.

\begin{quote}
\bn{আইনের} \bn{এই} \bn{ধারাটা} \bn{জটিল।}
\end{quote}

\begin{quote}
\emph{āiner ei dhārāṭā jôṭil.}
\end{quote}

\begin{quote}
This section of the law is complicated.
\end{quote}

\begin{cousinweb}
\textbf{\bn{ধারা}} (\emph{dhārā}) is Sanskrit \emph{dhārā}, a stream or flow (Monier-Williams). The sense 'a section of a law', one strand in the flow of the text, is a later legal use.
\end{cousinweb}

\begin{grammarlens}[title={What you've built}]
One more word for giving reasons and opinions.
\end{grammarlens}

\subsection*{Your turn}
\begin{itemize}
\raggedright
  \item \emph{Say it:} \emph{dhārā}
  \item \emph{Say it:} \emph{dhārā}, once more
  \item \emph{Say it:} \emph{āiner ei dhārāṭā jôṭil.} — This section of the law is complicated.
  \item \emph{Recall:} say the Bengali for lawful, then the Bengali for banned, then say \emph{dhārā} again
\end{itemize}

\begin{tcolorbox}[breakable,colback=teal!4,colframe=teal!35!black,title={Before you move on}]
What does \bn{ধারা} mean? (\textquotedblleft{}A section (of a law); a clause\textquotedblright{}.) Say it once more.
\end{tcolorbox}
\section[\texorpdfstring{\bn{পালন} \bn{করা} (pālôn kôrā)}{পালন করা (pālôn kôrā)}]{\bn{পালন} \bn{করা} (pālôn kôrā) — to observe (a rule, a duty); to carry out}

\label{lesson:BN-C340-palon-kora}



\begin{quote}
Before the new one: say the Bengali for banned, then the Bengali for a clause.
\end{quote}

\subsection*{You'll want to know: \bn{পালন} \bn{করা}}
\textbf{\bn{পালন} \bn{করা}} — \emph{pālôn kôrā} — \textquotedblleft{}to observe (a rule, a duty); to carry out\textquotedblright{}.

\begin{quote}
\bn{আমরা} \bn{সব} \bn{নিয়ম} \bn{পালন} \bn{করি।}
\end{quote}

\begin{quote}
\emph{āmrā sôb niyôm pālôn kôri.}
\end{quote}

\begin{quote}
We follow all the rules.
\end{quote}

Past (he/she …): \bn{তিনি} \bn{তাঁর} \bn{কর্তব্য} \bn{পালন} \bn{করলেন}, \emph{tini tā̃r kôrtôbyô pālôn kôrlen}.

\begin{cousinweb}
\textbf{\bn{পালন} \bn{করা}} (\emph{pālôn kôrā}) is Sanskrit \emph{pālana}, guarding and keeping (Monier-Williams), with \bn{করা}: to keep a rule, a duty or a festival day.
\end{cousinweb}

\begin{grammarlens}[title={What you've built}]
One more word for giving reasons and opinions.
\end{grammarlens}

\subsection*{Your turn}
\begin{itemize}
\raggedright
  \item \emph{Say it:} \emph{pālôn kôrā}
  \item \emph{Say it:} \emph{pālôn kôrā}, once more
  \item \emph{Say it:} \emph{āmrā sôb niyôm pālôn kôri.} — We follow all the rules.
  \item \emph{Recall:} say the Bengali for banned, then the Bengali for a clause, then say \emph{pālôn kôrā} again
\end{itemize}

\begin{tcolorbox}[breakable,colback=teal!4,colframe=teal!35!black,title={Before you move on}]
What does \bn{পালন} \bn{করা} mean? (\textquotedblleft{}To observe (a rule, a duty); to carry out\textquotedblright{}.) Say it once more.
\end{tcolorbox}
\section[\texorpdfstring{\bn{প্রয়োগ} \bn{করা} (prôyog kôrā)}{প্রয়োগ করা (prôyog kôrā)}]{\bn{প্রয়োগ} \bn{করা} (prôyog kôrā) — to apply; to enforce}

\label{lesson:BN-C340-proyog-kora}



\begin{quote}
Before the new one: say the Bengali for a clause, then the Bengali for to carry out.
\end{quote}

\subsection*{You'll want to know: \bn{প্রয়োগ} \bn{করা}}
\textbf{\bn{প্রয়োগ} \bn{করা}} — \emph{prôyog kôrā} — \textquotedblleft{}to apply; to enforce\textquotedblright{}.

\begin{quote}
\bn{পুলিশ} \bn{নতুন} \bn{আইনটা} \bn{প্রয়োগ} \bn{করছে।}
\end{quote}

\begin{quote}
\emph{pulish nôtun āinṭā prôyog kôrchhe.}
\end{quote}

\begin{quote}
The police are enforcing the new law.
\end{quote}

Past (he/she …): \bn{তিনি} \bn{পদ্ধতিটা} \bn{প্রয়োগ} \bn{করলেন}, \emph{tini pôddhôtiṭā prôyog kôrlen}.

\begin{cousinweb}
\textbf{\bn{প্রয়োগ} \bn{করা}} (\emph{prôyog kôrā}) is Sanskrit \emph{prayoga}, application or use (Monier-Williams), with \bn{করা}: to put a law, a method or a force to work.
\end{cousinweb}

\begin{grammarlens}[title={What you've built}]
One more word for giving reasons and opinions.
\end{grammarlens}

\subsection*{Your turn}
\begin{itemize}
\raggedright
  \item \emph{Say it:} \emph{prôyog kôrā}
  \item \emph{Say it:} \emph{prôyog kôrā}, once more
  \item \emph{Say it:} \emph{pulish nôtun āinṭā prôyog kôrchhe.} — The police are enforcing the new law.
  \item \emph{Recall:} say the Bengali for a clause, then the Bengali for to carry out, then say \emph{prôyog kôrā} again
\end{itemize}

\begin{tcolorbox}[breakable,colback=teal!4,colframe=teal!35!black,title={Before you move on}]
What does \bn{প্রয়োগ} \bn{করা} mean? (\textquotedblleft{}To apply; to enforce\textquotedblright{}.) Say it once more.
\end{tcolorbox}
